\documentclass[11pt]{article}
\usepackage[a4paper,margin=24mm]{geometry}
\usepackage{amsmath,amssymb,hyperref}
\title{A singleton counterexample to conjecture 00000008471}
\author{ziangni-sys}
\date{4 October 2026}
\begin{document}
\maketitle
The final clause asserts that the complement of the locked continuous
observables is a dense open set. We disprove density under the explicit
locking definition given in both versions of the conjecture.

Let $X=\{a\}$, with its unique compact topology, and let $T=\mathrm{id}_X$.
Its only Borel probability measure is $\delta_a$: any probability measure
assigns mass $1$ to $X=\{a\}$, and these are all its measurable sets.
This measure is invariant under $T$. For every continuous
$f:X\to\mathbb R$ we have
\[
\int_X f\,d\delta_a=f(a).
\]
Since there is no other invariant probability measure, the actual
maximizing-measure set is
\[
M(f)=\{\delta_a\}\qquad\text{for every }f\in C(X,\mathbb R).
\]
This conclusion follows from maximization over all invariant probability
measures, rather than from a restricted list of candidates.

The point $a$ has period $1$, and the equidistribution on its orbit is
\[
\frac1{1}\sum_{j=0}^{0}\delta_{T^j(a)}=\delta_a.
\]
Consequently every continuous observable is locked in the stated sense:
its entire maximizing-measure set consists of the equidistribution on
one periodic orbit. In particular the locked-function set is the whole
space $C(X,\mathbb R)$. Its complement is empty. The function space is
nonempty (it contains the zero function), so the empty set is not dense
in its usual uniform topology. Thus the claimed dense open complement
law is false. Although the empty complement is open, it fails the
required density.

Neither source excludes singleton dynamical systems. Even if locking
also requires persistence under perturbation, this example satisfies
that condition: every continuous perturbation still has the same unique
maximizing measure. The counterexample concerns the final conjunct;
it does not assert a failure of the generic-locking conjunct.

\paragraph{Formal verification.}
The accompanying Lean project uses the actual compact space
\texttt{Unit}, its Borel probability measures, the actual identity
map, Bochner integrals and \texttt{ContinuousMap Unit Real}. It proves
uniqueness of every probability measure, the full maximizing-measure
characterization, the weighted orbit measure and periodicity, locking
of every observable, and failure of Mathlib's actual \texttt{Dense}
predicate for the complement. The project is pinned to Lean 4.19.0 and
Mathlib commit \texttt{c44e0c8ee63ca166450922a373c7409c5d26b00b}.
\end{document}
